\documentclass[11pt]{article}
% === core_packages.tex ===
\usepackage{amsmath, amssymb, amsthm, amsfonts}
\usepackage{geometry}
\usepackage{hyperref}
\usepackage{bookmark}
\usepackage{graphicx}
\usepackage{physics}
\usepackage{booktabs}
\usepackage{listings}
\usepackage{amscd}
\usepackage{tikz-cd}
\usepackage{mathtools}
\usepackage{xspace}
\usepackage[backend=biber, style=numeric]{biblatex}
\usepackage{tcolorbox}
\tcbuselibrary{theorems}
\usepackage{parskip}
\usepackage{tikz}
\usepackage{etoolbox}
\usetikzlibrary{arrows.meta, shapes, positioning, calc, cd, shapes.geometric, backgrounds}
\usepackage[en-US]{datetime2}
\usepackage{array}

% --- Theory names ---
\newcommand{\iame}{[theory needs name]}
\newcommand{\iameverything}{[theory needs name]}
\newcommand{\theory}{[theory needs name]}


% --- Math & Ontology Macros ---
\newcommand{\entropy}{S}
\newcommand{\opponent}{opponent}
\newcommand{\aside}[1]{[\emph{\opponent}: ``#1'']}
\newcommand{\paper}[1]{[\emph{#1}: ``#1'']}
\newcommand{\fabric}{voxel}
\newcommand{\Fabric}{Voxel}
\newcommand{\pdflink}[1]{}



% --- Global Layout Defaults ---
\geometry{letterpaper, margin=1in}
\hypersetup{
    colorlinks=true,
    linkcolor=black,
    citecolor=black,
    filecolor=magenta,
    urlcolor=blue
}
\hfuzz=5pt \vfuzz=5pt \hbadness=10000 \vbadness=10000

\author{\iame}


\usepackage{tocloft}

% --- Article Specific Theorems ---
\newtheorem{lemma}{Lemma}
\theoremstyle{definition}
\newtheorem{axiom}{Axiom}
\newtheorem{theorem}{Theorem}
\newtheorem{corollary}{Corollary}
\newtheorem{definition}{Definition}
\newtheorem{proposition}{Proposition}
\newtheorem{conjecture}{Conjecture}

\makeatletter
\let\remark\@undefined
\let\endremark\@undefined
\makeatother
\theoremstyle{remark}
\newtheorem*{remark}{Remark}

\newtcbtheorem[number within=section]{principle}{Principle}%
{colback=blue!5,colframe=blue!50!black,fonttitle=\bfseries}{principle}

% --- Article TOC Layout ---
\setlength{\cftbeforesecskip}{1.0em}
\renewcommand{\cftsecleader}{\cftdotfill{\cftdotsep}}
\renewcommand{\cftsubsecleader}{\cftdotfill{\cftdotsep}}


\addbibresource{../references.bib}

\newcommand{\unarydigit}{\circ}      % a unary digit
\newcommand{\bnd}{\,\big|\,}         % boundary between marked and unmarked

\title{From Unary to Binary System\\
\large Observer-relative distinction without a binary primitive}
\author{Juha Meskanen}
\date{2019}

\begin{document}
\maketitle

\begin{abstract}
Is binary an essential primitive of information? We argue that it is not.
Dividing a featureless infinity yields parts that are interchangeable and carry
no intrinsic property: unary digits. We make one further assumption, that an
observer exists. There is no outside, so the observer is all of $U$, and what it
experiences is a traversal through readings of $U$. Time is the order of the
traversal, and binary distinction is the relation between successive frames; it is
not a property of the digits. A successor traversal yields the natural numbers.
A traversal of the subsets of $n$ digits yields the binary counter, with the
significance of each digit fixed by its role in the traversal and not by the
digit. Information is therefore emergent, a property of the observer, and binary
is one particularly compact alphabet for it. The scope is the transition from
unary to binary distinction. Which traversals contain observers is a question of
measure and is not treated here.
\end{abstract}

\section{Initial assumptions}

\begin{itemize}
  \item The deep nature of everything is infinite. Were it finite we would have to
        ask why $N$ and not $N+1$. ``Infinite'' means not finite; nothing below
        depends on which infinity.
  \item It can be anything, and it is everything. Otherwise we would have to ask
        why some things are forbidden and others not.
  \item It has no privileged structure, no distinguished state, no intrinsic
        information, mathematics, geometry or physical law. Otherwise we would have
        to ask where they came from.
  \item All ways of introducing distinctions are equally unprivileged, so all of them exist.
  \item There is no external metaphysics, signal or agent performing operations.
        Divisions, readings and traversals below are relations that exist, not acts.
  \item The nature of everything is abstract. Otherwise we would have to ask where
        the substrate came from.
\end{itemize}
A single further assumption, the existence of an observer, is introduced in
Section~\ref{sec:observer}.

\section{From infinity to unary digits}\label{sec:unary}

Let $X$ be the featureless infinity. Consider any division of $X$ into parts.
By invariance under differentiation, dividing a structureless infinity cannot
introduce intrinsic properties into the parts. Each part is again a featureless
infinity of the same kind as $X$. The parts are therefore interchangeable and carry
no intrinsic property, size included. A part has exactly one state. We call it a
\emph{unary digit}.

A division into $n$ parts gives a \emph{unary collection} $U_n$. The number $n$
counts parts and says nothing about their content, which is the same infinity in
each. It is the only invariant of the collection, since $U_n$ is invariant under
$\mathrm{Sym}(n)$: no ordering, no meaning, no distinction between digits is
intrinsically defined.

Divisions into every number of parts exist, by the fourth assumption. A finite
$n$ is a property of a division and not of $X$, so ``why this $n$'' needs no
answer. We write $U$ for the collection under consideration when $n$ does not
matter.

\section{The observer, readings and time}\label{sec:observer}

\textbf{Assumption.} An observer exists.

$X$ has no outside, so the observer cannot be external to $U$. We take it to be
all of $U$. Digits the observer does not traverse are not observed and are
irrelevant to it; asking about them is asking about a universe outside the universe.

\paragraph{Readings.} As a bare collection, $U_n$ has a single arrangement, since
$\mathrm{Sym}(n)$ identifies them all. Distinct arrangements exist only as
\emph{readings}: assignments of roles to the digits, for instance which are
marked. A reading is not a property of the digits. It is held on the reader's
side, as the field layout of a 64-bit floating-point word is: no separator is
stored in the bits, and two interpreters with different conventions read the same
bits as different values.

No attribute of $U$ separates readings that exist from readings that do not.
$\mathrm{Sym}(n)$ maps every reading onto others that nothing in $U$
distinguishes from it, so the whole orbit exists, and with the fourth assumption,
all readings exist.

\paragraph{Traversal and time.} A \emph{frame} is a reading. A \emph{traversal} is
an ordered sequence of frames, and the observer is the traversal. Time is its
order. An observer whose traversal is a single frame has no before and after and
no time. Two frames differ only as readings, so their difference is relational and
needs no intrinsic state of the digits. Most traversals are noise, and which ones
carry an observer is a question of measure that this paper does not treat.

\paragraph{Where the binary distinction enters.} Two-ness enters in two places,
and neither is an intrinsic two-valued state of a digit. Within a frame, a digit
is marked or not by the reader's convention. Between frames, the distinction is
before versus after. The smallest case is $n=1$: a single digit and two frames,
$\varnothing$ and $\{\unarydigit\}$. It needs a traversal of length two and nothing
more.

\section{Natural numbers from a traversal}\label{sec:numbers}

Take a traversal of nested readings $S_0\subset S_1\subset\dots\subset S_n$ with
$|S_k|=k$, where $S_k$ is the set of marked digits of frame $k$ and the split of the
frame is $\unarydigit^k\bnd\unarydigit^{n-k}$:
\[
\begin{array}{c|c}
\text{Frame} & \text{Split}\\\hline
0 & \bnd\unarydigit\unarydigit\unarydigit\\
1 & \unarydigit\bnd\unarydigit\unarydigit\\
2 & \unarydigit\unarydigit\bnd\unarydigit\\
3 & \unarydigit\unarydigit\unarydigit\bnd
\end{array}
\]
The natural number $k$ is the position of the frame in the traversal, and the
passage from $k$ to $k+1$ is the successor. No counting outside the observer is
required: the observer's own progress is the count. Which digit is added at each step
is not fixed by $U$. The traversal individuates the digits by their order of
addition, and the digits themselves acquire no identity.

Zero is the initial frame. For infinite $n$ the traversal continues past every
finite stage, and we do not develop this. The interpretation is the observer's, as
usual: a bare unary collection has no arithmetic. An implementation of the
resulting operations on strings of marks is given in the supplementary script
\texttt{unary.py}.

\section{Binary from a traversal}\label{sec:binary}

Now let the observer traverse all $2^n$ subset readings of $U_n$ in the order of a
binary counter: in frame $k$, digit $j$ is marked iff bit $j$ of $k$ is $1$. The
digits carry no labels. Nevertheless digit $j$ is identified by its flip period
$2^j$, which is a property of the traversal and not of the digit. If the digits
are relabelled by any $\pi\in\mathrm{Sym}(n)$, the raw frames change but the
set of flip periods does not. The positional weights of
\[
N=\sum_{j=0}^{n-1} b_j\,2^j
\]
are thus roles in the traversal, and a binary numeral is the reading of a single frame
of this traversal. In the script \texttt{traversal.py} we check this for $n=8$:
256 distinct frames, roles $1,2,4,\dots,128$ recovered from flip timing alone,
and the same under random relabelling.

\paragraph{Comparison with the unary chain.} The nested traversal of
Section~\ref{sec:numbers} reaches $2^n$ frames only with $2^n-1$ digits. The
binary counter reaches them with $n$ digits, at the price of requiring all
subset readings, which exist by Section~\ref{sec:observer}. The $n$-bit
string is thus the compressed address of a position in the unary traversal;
conversion between the two is pairing off marks (unary to binary) and doubling with
carry (binary to unary). Nothing about this is intrinsic to the digits.

\paragraph{Positional weights.} Positional weighting is additional structure
supplied by the reader. The unary case $w=1$ is its degenerate form,
$N_1=\sum_j b_j$, which is independent of order for any reader. This is not
evidence for the symmetry of $U$: any symmetric function of the marks has that property.

\section{Consequences}

\[
\begin{aligned}
&\text{Featureless infinity}\\
&\quad\downarrow\ \text{division}\\
&\text{Unary digits (interchangeable, no intrinsic property)}\\
&\quad\downarrow\ \text{observer: all of } U\text{; readings exist}\\
&\text{Traversal of frames: time, before/after}\\
&\quad\downarrow\\
&\text{Successor counter; binary counter (roles by flip period)}.
\end{aligned}
\]
Information is not assumed to exist in $U$. Distinction arises relative to a
traversal, and binary is the simplest alphabet in which a traversal can be
addressed compactly.

\section{Assumptions}

\begin{enumerate}
  \item An observer exists. Since $X$ has no outside, it is all of $U$.
  \item All readings of $U$ exist. This follows from the fourth assumption of
        Section~1 and from the absence of any attribute that could select among
        them. We state it separately for transparency.
  \item Marking, order of addition, positional weight and numeral labels are
        conventions on the reader's side. None is a property of the digits.
  \item Which traversals carry observers, that is, the measure, is not treated.
  \item Natural numbers require finite $n$. The infinite case is not developed.
  \item (Optional) The observer is finite: $n<\infty$. This is motivated by the
        empirical observation that observers do not encounter actual infinities.
\end{enumerate}

\section{Conclusion}

A featureless infinity, divided, gives indistinguishable unary digits. An observer
that is all of them, and therefore has no outside, can have an experience only as a
traversal through their readings. Time, number and binary representation are then
properties of the traversal, and the digits stay unary throughout.

\medskip
\noindent\textbf{Existence does not require information. Information requires
distinction, and distinction can emerge internally through observation.}

Information, time and the laws of physics are emergent properties of the observer. The
universe at its deepest level is neither temporal nor informational, and from
inside it cannot be told from nothing.

\end{document}
